\section{Traditions Examined}
\label{sec:traditions}

Everything vivid in Thomas's biography comes from one archive: the
canonization dossier---the inquiries of 1319 and 1321 and the Lives of
William of Tocco (1318--23), Bernard Gui, and Peter Calo (c.~1300)---
assembled by devoted men, for the purpose of proving sanctity, forty to
ninety years after the events of the youth they describe. Such evidence
establishes memory, cult, and the shape of a reputation; it cannot, by
itself, backdate a private scene to 1244. This section examines the
memorable traditions one by one, without credulity and without ridicule.

\subsection{The temptress and the girdle}

The story: during the family detention his brothers introduced a woman
into his room to seduce the novice; Thomas snatched a burning brand from
the fire and drove her out, marking a cross on the door; asleep, he was
girded by angels with---in the wording the 1912 Encyclopedia
transmits---``the girdle of perpetual virginity.'' The chain: the
tradition itself says Thomas told no one but Reginald, and Reginald only
near the end; every witness is thus canonization-era, and the scene
reaches the record as a reported deathbed confidence. The status: the
detention is historical; the bedroom scene is unverifiable and its
literary shape---trial, victory, heavenly investiture---is the shape
hagiography loves. Its reception, however, is real history: the
``angelic warfare'' of the story made Thomas a patron of chastity, and
the girdle a devotional emblem. This study asserts the detention,
reports the story with its chain, and claims nothing about the relic
traditions attached to the girdle.

\subsection{The dumb ox}

The Cologne anecdote (section above) has the same profile: late
witnesses, no contemporary record, a prophecy remembered after its
spectacular fulfillment---the classic form of retrospective
master-recognizes-genius stories. It may preserve a real nickname and a
real judgment of Albert's; it cannot be verified; and its function in
the Lives is transparent, converting the saint's youthful silence into
foreordained greatness. This study's searches found no witness earlier
than the canonization generation (a bounded negative result).

\subsection{The straw remark}

The witness chain of 6 December 1273 is set out beside the event itself
(with the cessation secure and the wording second-hand); it belongs in
this audit as the tradition's most instructive case: a saying
universally quoted, resting on one sworn but hearsay deposition of
1319, whose English forms disagree because the Latin memory was itself
a report. The remark's fame is modern as well as medieval---the
Stanford Encyclopedia quotes it, Benedict~XVI retold it---and none of
that reception adds a witness.

\subsection{The speaking crucifix and the levitations}

The Naples story---the sacristan's vision of the friar raised in the
air before the crucifix in St.~Nicholas's chapel, and the voice: ``Thou
hast written well of me, Thomas''---is canonization-era memory, and the
1912 Encyclopedia honestly notes that ``similar declarations are said
to have been made at Orvieto and at Paris.'' The multiplication is the
audit's datum: the scene was how the tradition expressed its conviction
that the doctrine and the sanctity were one thing, the truth of the
books certified by the Christ they describe. Benedict~XVI retold the
story in 2010 as a spiritual exemplum; a papal catechesis hands on the
tradition and does not convert it into an eyewitness record. This study
narrates it as the cause's memory, invents no natural explanation, and
certifies no miracle.

\subsection{The restless corpse: relic history}

The posthumous history of Thomas's body is documented in layers of very
unequal strength.

\begin{itemize}
\item \textbf{Secure core.} Death and burial at Fossanova, 1274; the
Cistercians' jealous custody of so valuable a grave; and the translation
of the remains to the Dominicans of Toulouse, arriving 28 January 1369
(the 1912 Encyclopedia; the medieval custody conflicts and the 1369
ceremonies are the subject of R\"as\"anen's 2017 monograph, cited at
record level). The canonization-era sources already tell of exhumations
and translations within Fossanova, and later accounts of the Cistercian
custody include grim expedients---separation of the head, even boiling
of the remains---which this study reports only as tradition-history at
the reported level of the modern scholarship, not as verified events.
\item \textbf{Reported extensions.} During the French Revolution the
relics are reported to have been moved from the suppressed Jacobins
church to the basilica of Saint-Sernin, and in 1974, for the seventh
centenary, returned to the Jacobins, where the shrine now stands
(current public accounts checked 2026-07-25; no official custodial
document was checked, and no current authentication is asserted).
\item \textbf{Competing claims.} Head and arm relics are claimed
elsewhere---the checked modern account names San Domenico Maggiore in
Naples; other head traditions attach to Priverno, near Fossanova. This
study adjudicates none of them. The cult, it should be said plainly,
does not depend on the adjudication: the Church honors the saint, not a
forensically certified skeleton.
\end{itemize}

\subsection{Iconography}

The standard image---the corpulent friar with book, sun or star on the
breast, sometimes a dove at the ear, often a chained Averroes beneath
the feet---is reception, not portraiture: no verified life-portrait
exists, and the attributes encode the reputation (doctrine radiating
from the breast, inspiration, the defeated adversary). The ox and the
girdle entered art from the anecdotes audited above; the images are
evidence for the centuries that painted them, and this study uses them
for nothing else.
